% Bloom levels: remember, understand, apply, analyze, evaluate, create.
\begin{problema}\label{prob:d20c80f}
State, without calculation, the formula for the expected frequency $E_{ij}$ under the null hypothesis of independence in an $r\times c$ contingency table, and the formula for the corresponding degrees of freedom.
\end{problema}

\begin{problema}\label{prob:b002759}
Explain, without numerical data, why
\[
	E_{ij}=\frac{R_iC_j}{N}
\]
is the correct expected frequency under independence, starting from the fact that if two events $F_i$ and $C_j$ are independent, then $P(F_i\cap C_j)=P(F_i)P(C_j)$.
\end{problema}

\begin{problema}\label{prob:147e151}
A road-safety audit checks whether vehicle type and weekday are independent in urban accidents ($n=300$):
\[
\begin{array}{|l|c|c|c|c|}\hline
\text{Vehicle type} & \text{Monday} & \text{Wednesday} & \text{Friday} & \text{Total}\\\hline
\text{Car} & 30 & 25 & 45 & 100\\\hline
\text{Van} & 20 & 30 & 30 & 80\\\hline
\text{Motorcycle} & 25 & 20 & 35 & 80\\\hline
\text{Truck} & 15 & 15 & 10 & 40\\\hline
\text{Total} & 90 & 90 & 120 & 300\\\hline
\end{array}
\]
Compute the expected frequencies, the $\chi^2$ statistic, the degrees of freedom, and decide at $\alpha=0.05$, using $\chi^2_{0.05,6}=12.59$, whether there is a significant relationship between day and vehicle type.
\end{problema}

\begin{problema}\label{prob:3c6b064}
For a $2\times2$ contingency table with cells $a,b,c,d$, row totals $R_1,R_2$, and column totals $C_1,C_2$, prove algebraically that Pearson's general statistic under independence collapses to
\[
	\chi^2=\frac{N(ad-bc)^2}{(a+b)(c+d)(a+c)(b+d)}.
\]
Explain the meaning of $(ad-bc)$ as the determinant of the contingency matrix, relating it to the odds ratio $\operatorname{OR}=ad/bc$.
\end{problema}

\begin{problema}\label{prob:878a46f}
After rejecting $H_0$ in a chi-square independence test between ``advertising channel seen'' and ``purchase decision,'' a marketing analyst concludes: ``it has been proved that the advertising channel causes the purchase decision.'' Evaluate this conclusion: what exactly does a significant $\chi^2$ independence test prove, and what does it not prove? Propose an alternative explanation, such as a confounding variable, that could produce the same non-independence without a direct causal relationship.
\end{problema}

\begin{problema}\label{prob:a5524c8}
Design an original example, with context and an $r\times c$ contingency table with $r\ge2$ and $c\ge2$, where a $\chi^2$ independence test is applied. Use a context different from the traffic-accident example in this section. Compute at least one expected cell frequency and state the degrees of freedom of your table.
\end{problema}

\begin{solproblema}[prob:d20c80f]
\[
	E_{ij}=\frac{R_iC_j}{N},
\]
where $R_i$ is the row total, $C_j$ is the column total, and $N$ is the grand total. The degrees of freedom are
\[
	\nu=(r-1)(c-1).
\]
\end{solproblema}

\begin{solproblema}[prob:b002759]
Under independence, the marginal probabilities are estimated by $P(F_i)\approx R_i/N$ and $P(C_j)\approx C_j/N$. Thus
\[
	P(F_i\cap C_j)=P(F_i)P(C_j)\approx\frac{R_i}{N}\frac{C_j}{N}.
\]
Multiplying by $N$ converts this probability into an expected count:
\[
	E_{ij}=N\frac{R_i}{N}\frac{C_j}{N}=\frac{R_iC_j}{N}.
\]
\end{solproblema}

\begin{solproblema}[prob:147e151]
The expected frequencies are: Car $(30,30,40)$; Van $(24,24,32)$; Motorcycle $(24,24,32)$; Truck $(12,12,16)$. The statistic is approximately
\[
	\chi^2=8.49.
\]
The degrees of freedom are
\[
	\nu=(4-1)(3-1)=6.
\]
Since $8.49<12.59$, do not reject $H_0$: there is no statistically significant relationship between day and vehicle type at the $5\%$ level.
\end{solproblema}

\begin{solproblema}[prob:3c6b064]
For example,
\[
	E_{11}=\frac{(a+b)(a+c)}{N},
\]
and simplifying $O_{11}-E_{11}$ gives a term proportional to $ad-bc$. The four cell deviations have the same determinant-driven magnitude with alternating signs. Substituting these deviations into
\[
	\chi^2=\sum\frac{(O_{ij}-E_{ij})^2}{E_{ij}}
\]
and simplifying yields
\[
	\chi^2=\frac{N(ad-bc)^2}{(a+b)(c+d)(a+c)(b+d)}.
\]
The term $ad-bc$ is the determinant of the table matrix. If $ad=bc$, the determinant is zero, the rows are proportional, $\operatorname{OR}=ad/bc=1$, and the table is consistent with independence.
\end{solproblema}

\begin{solproblema}[prob:878a46f]
The conclusion is incorrect. A significant independence test shows statistical association between the two categorical variables; it does not establish causality or direction of effect. A plausible confounder is consumer income: income could affect both which advertising channels consumers see and their probability of purchase, producing association between channel and purchase even if the channel itself is not the direct cause.
\end{solproblema}

\begin{solproblema}[prob:a5524c8]
Example: a university studies whether study modality and course result are independent among $400$ students:
\[
\begin{array}{|l|c|c|c|}\hline
\text{Modality} & \text{Passed} & \text{Failed} & \text{Total}\\\hline
\text{In person} & 120 & 30 & 150\\\hline
\text{Hybrid} & 90 & 30 & 120\\\hline
\text{Online} & 80 & 50 & 130\\\hline
\text{Total} & 290 & 110 & 400\\\hline
\end{array}
\]
For the In person--Passed cell,
\[
	E_{11}=\frac{150(290)}{400}=108.75.
\]
The degrees of freedom are $(3-1)(2-1)=2$.
\end{solproblema}
